% !PRINTED 155
\weberpage{155}
Just as easily we obtain:

\par
5.\enspace\gesperrt{By composition with the identical substitution $J$ no substitution is changed.}

The substitution $J$ is therefore regarded in composition as a \gesperrt{unit}, and, wherever it occurs compounded with other substitutions, it may be omitted.

Now the following theorem holds further:

\par
6.\enspace\gesperrt{To every substitution $A$ there is one and only one inverse substitution $A^{-1}$, by which the condition}

\weq{15}{AA^{-1} = A^{-1}A = J}

\gesperrt{is satisfied.}

This last theorem follows from the fundamental formulae of the theory of determinants, if one determines the elements $\alpha_h^{(k)}$ of $A^{-1}$ from the linear equations

\weq{16}{\begin{aligned} \sum^i a_i^{(h)}\alpha_k^{(i)} &{}= 0, & h \gtrless k\\ &{}= 1, & h = k \end{aligned}}

from which, if we denote by $A_h^{(k)}$ the subdeterminants of $|A|$, there results

\weq{17}{|A|\,\alpha_h^{(k)} = A_k^{(h)}}

and which entail the other system

\weq{18}{\begin{aligned} \sum^i \alpha_i^{(h)} a_k^{(i)} &{}= 0, & h \gtrless k\\ &{}= 1, & h = k \end{aligned}}

The inverse substitution to $A$ is nothing other than the solution of the system of equations (1) of the direct substitution $A$, so that from $(y) = A(x)$ it follows that:

\weq{19}{(x) = A^{-1}(y).}

For the composition of the inverse substitutions there results from $ABB^{-1}A^{-1} = J$ the theorem:

\weq{20}{(AB)^{-1} = B^{-1}A^{-1}.}

If $A$, $B$, $C$ are three substitutions, then by composition with $A^{-1}$ it follows from each of the two equations

\[
AB = AC, \quad BA = CA,
\]

that $B = C$ must hold. Accordingly, for the composition of the substitutions the characteristic marks of a group are fulfilled, and thus the totality of all substitutions is
%CONT>>
